% © 2026 Ing. David Jaroš — CC BY-NC-ND 4.0
%
% This work is licensed under a Creative Commons Attribution-NonCommercial-NoDerivatives
% 4.0 International License (CC BY-NC-ND 4.0).
%
% License History: Earlier drafts (up to v0.3) were released under CC BY 4.0.
% From v0.4 onward, all material is released under CC BY-NC-ND 4.0 to protect
% the integrity of the theoretical work during ongoing academic development.

\documentclass[11pt,a4paper]{article}
\usepackage{amsmath, amssymb, amsthm}
\usepackage{hyperref}
\usepackage{geometry}
\geometry{margin=2.5cm}

\newtheorem{theorem}{Theorem}
\newtheorem{proposition}{Proposition}
\newtheorem{definition}{Definition}
\newtheorem{remark}{Remark}
\newtheorem{conjecture}{Conjecture}

\title{Dynamical Equations of the \texorpdfstring{$\Theta$}{Theta}-Field in UBT\\
\large Field Equations, Euler--Lagrange Structure, and Particle Spectra}
\author{Ing.\ David Jaroš}
\date{2026}

\begin{document}
\maketitle

\begin{abstract}
We write the full dynamical field equations derived from the biquaternionic
$\Theta$-field action in Unified Biquaternion Theory (UBT) and analyse how
particle spectra emerge.  The Euler--Lagrange equations are derived
component-by-component, the KK decomposition is used to obtain the 4D
effective equations for each mode, and the linearised spectrum around the
vacuum is computed.  All claims are labelled with their proof status.
\end{abstract}

\tableofcontents

% ======================================================================
\section{Euler--Lagrange Equations}
\label{sec:el}
% ======================================================================

\subsection{Variation of the action}

Starting from the canonical action (eq.\ (1) of
\texttt{research/theta\_quantum\_field.tex} and
\texttt{consolidation\_project/appendix\_A\_theta\_action.tex}):
\begin{equation}
\label{eq:action}
  S[\Theta]
  = \int d^4x\, d\tau\;\sqrt{|\det G|}\;
    \Bigl[\tfrac{1}{2}G^{MN}\langle D_M\Theta, D_N\Theta\rangle
          - V(\Theta)
          - \tfrac{1}{4}\langle F_{MN}, F^{MN}\rangle\Bigr],
\end{equation}
the Euler--Lagrange equations follow from $\delta S / \delta\Theta^\dagger = 0$.

\begin{theorem}[Θ-field equation of motion]\label{thm:eom}
  The equation of motion for $\Theta$ derived from \eqref{eq:action} is:
  \begin{equation}
  \label{eq:eom}
    -D^M D_M\,\Theta + V'(\Theta) = 0,
  \end{equation}
  or equivalently in terms of the biquaternionic d'Alembertian:
  \begin{equation}
  \label{eq:box}
    \Box_{\mathbb{B}}\,\Theta \;=\; V'(\Theta),
  \end{equation}
  where $\Box_{\mathbb{B}} = G^{MN}D_M D_N$ is the biquaternionic
  covariant d'Alembertian and $V'(\Theta) = \partial V/\partial\Theta^\dagger$.
  \textbf{Status: [L1]} — proved given action \eqref{eq:action} and axioms.
\end{theorem}

\begin{remark}[Relation to master equation]
  Equation \eqref{eq:eom} is the Euler--Lagrange form of the UBT master
  equation $\nabla^\dagger\nabla\Theta = \kappa\mathcal{T}$ in the absence
  of external currents.  The two forms agree when $V'(\Theta) = \kappa\mathcal{T}$.
  Source: \texttt{canonical/fields/theta\_field.tex} §Field Equations.
\end{remark}

% ======================================================================
\section{Component Field Equations}
\label{sec:components}
% ======================================================================

\subsection{Quaternionic components}

Expanding $\Theta = \Theta_0 + \Theta_1 i + \Theta_2 j + \Theta_3 k$
(quaternion basis), \eqref{eq:eom} splits into four coupled equations
for the complex scalars $\Theta_a$, $a = 0,1,2,3$:
\begin{equation}
\label{eq:component_eom}
  -G^{MN}D_M D_N\,\Theta_a + \sum_{b=0}^{3} M_{ab}^2\,\Theta_b
  = J_a,
\end{equation}
where $M_{ab}^2 = \partial^2 V/(\partial\Theta_a\partial\Theta_b^\dagger)|_{\Theta=v}$
is the mass matrix evaluated at the vacuum $v$, and $J_a$ are source currents.

\begin{remark}[Status]
  Equation \eqref{eq:component_eom} is a set of four coupled
  Klein--Gordon-type equations with a $4\times 4$ mass matrix.  The
  diagonalisation of $M_{ab}^2$ produces the physical particle spectrum.
  Status: \textbf{[L1]}.
\end{remark}

% ======================================================================
\section{4D Effective Equations from KK Reduction}
\label{sec:kk_eff}
% ======================================================================

\subsection{Mode-by-mode equations}

Inserting the KK expansion $\Theta_a(x,t,\psi) = \sum_n \hat\Theta_{a,n}(x,t)\,
e^{in\psi/R_\psi}$ into \eqref{eq:component_eom}:

\begin{proposition}[4D effective equation for KK mode $n$]\label{prop:eff_eq}
  \begin{equation}
  \label{eq:kk_eff}
    \Bigl(-\partial^\mu\partial_\mu + m_{a,n}^2\Bigr)\hat\Theta_{a,n}(x,t)
    = \hat J_{a,n}(x,t),
  \end{equation}
  where $m_{a,n}^2 = M_{aa}^2 + n^2/R_\psi^2$ (no sum on $a$) for the
  diagonal sector, and $\hat J_{a,n}$ collects interaction vertices.
  \textbf{Status: [L1]}.
\end{proposition}

\subsection{Low-energy spectrum ($|n|/R_\psi \ll M_{aa}$)}

In the low-energy limit where all KK modes with $n \neq 0$ decouple
($R_\psi^{-1} \gg M_{aa}$), the effective 4D theory contains only the
zero-mode $n=0$ fields:
\begin{itemize}
  \item $\hat\Theta_{0,0}$ — massless scalar (Higgs zero mode)
  \item $\hat\Theta_{1,0}, \hat\Theta_{2,0}, \hat\Theta_{3,0}$ — vector fields
        (gauge bosons in low-energy limit)
\end{itemize}

\begin{remark}[Higgs and gauge bosons as zero modes]
  The identification of $n=0$ modes with Higgs and gauge fields follows
  the standard Kaluza--Klein argument.  It is consistent with the gauge
  group derivation in \texttt{canonical/bridges/gauge\_emergence\_bridge.tex}.
  Status: \textbf{[L1]}.
\end{remark}

% ======================================================================
\section{Mass Generation}
\label{sec:mass}
% ======================================================================

\subsection{Vacuum and symmetry breaking}

The potential is taken as the Mexican-hat form:
\begin{equation}
\label{eq:potential}
  V(\Theta) = \tfrac{1}{4}\lambda\bigl(|\Theta|^2 - v^2\bigr)^2,
\end{equation}
with $v > 0$ real.  The vacuum manifold is
$\{|\Theta| = v\} \cong S^7 \subset \mathbb{H}\otimes\mathbb{C}$
(a 7-sphere in the biquaternion algebra).

\begin{proposition}[Mass matrix from symmetry breaking]\label{prop:mass_matrix}
  Expanding around the vacuum $\Theta = v + h + i\phi_a\, e_a$
  ($h$ Higgs, $\phi_a$ Goldstone), the mass matrix at tree level is:
  \begin{equation}
  \label{eq:mass_matrix}
    M_h^2 = 2\lambda v^2, \qquad M_{\phi_a}^2 = 0 \quad (a=1,2,3).
  \end{equation}
  The three Goldstone modes $\phi_a$ are eaten by the gauge bosons via the
  Higgs mechanism.
  \textbf{Status: [L1]}.
\end{proposition}

\subsection{Fermion masses from Yukawa coupling}

From \texttt{research/theta\_fermion\_emergence.tex} §4, the Yukawa
Lagrangian is $\mathcal{L}_Y = -y\bar\psi\phi\chi + \mathrm{h.c.}$
After Higgs condensation $\langle\phi\rangle = v$:
\begin{equation}
\label{eq:fermion_mass}
  m_f = y\,v.
\end{equation}

\begin{remark}[Status and open problem]
  The relation \eqref{eq:fermion_mass} is standard and Status: \textbf{[L1]}.
  However, the Yukawa coupling $y$ and the vacuum value $v$ are not yet
  derived from $R_\psi$ and the moduli.  See \texttt{research/fermion\_mass\_generation.tex}
  for the detailed status of this open problem.
\end{remark}

% ======================================================================
\section{KK Mode Effective Mass Matrix}
\label{sec:kk_mass_matrix}
% ======================================================================

\subsection{Coupling between KK modes}

The potential $V(\Theta)$ mixes different KK modes.  For the cubic term
$V_3 = \lambda v |\Theta|^3$, the coupling between modes $n_1, n_2, n_3$
is:
\begin{equation}
\label{eq:kk_coupling}
  C_{n_1 n_2 n_3}
  = \lambda v \int_0^{2\pi R_\psi} \!\! d\psi\;
    e^{i(n_1+n_2+n_3)\psi/R_\psi}
  = 2\pi R_\psi\, \lambda v\, \delta_{n_1+n_2+n_3,0}.
\end{equation}

KK-number conservation is exact at tree level (Status: \textbf{[L0]}).

\subsection{Effective mass matrix in mode space}

Diagonalising the full mass matrix including KK mixing, the physical masses
of the lightest modes (quartic potential) are:
\begin{equation}
\label{eq:kk_mass_eff}
  m_{n,\mathrm{phys}}^2
  = M_h^2 + \frac{n^2}{R_\psi^2}
    + \mathcal{O}\!\Bigl(\frac{\lambda v^2}{R_\psi^2}\Bigr).
\end{equation}

\begin{remark}[Status]
  Equation \eqref{eq:kk_mass_eff} is the leading-order result.
  Higher-order KK mixing corrections are suppressed by $\lambda v^2 R_\psi^2$.
  Status: \textbf{[L1]}.
  Comparison with the Standard Model mass spectrum (Higgs, $W$, $Z$, fermions)
  is discussed in \texttt{research/fermion\_mass\_generation.tex}.
\end{remark}

% ======================================================================
\section{Propagation in Curved Background}
\label{sec:curved}
% ======================================================================

In the presence of gravity (non-flat $g_{\mu\nu}$), the d'Alembertian
$\Box_{\mathbb{B}}$ in \eqref{eq:box} receives curvature corrections:
\begin{equation}
  \Box_{\mathbb{B}}\,\Theta
  = \frac{1}{\sqrt{|g|}}\partial_M(\sqrt{|g|}\,G^{MN}\partial_N\Theta)
    + \xi R\,\Theta + \ldots,
\end{equation}
where $\xi$ is the non-minimal coupling to curvature and $R$ is the Ricci
scalar.  The minimal coupling choice $\xi = 0$ is assumed unless stated
otherwise.

\begin{remark}[GR compatibility]
  As established in \texttt{canonical/bridges/GR\_chain\_bridge.tex},
  UBT fully embeds General Relativity.  The Θ-field propagates on the
  emergent metric $g_{\mu\nu} = \mathrm{Re}\,\mathrm{Tr}(\partial_\mu\Theta\,\partial_\nu\Theta^\dagger)$.
  Status: \textbf{[L1]}.
\end{remark}

% ======================================================================
\section{Open Problems}
\label{sec:open}
% ======================================================================

\begin{enumerate}
  \item \textbf{[O] Non-minimal coupling $\xi$}: Whether $\xi = 0$ or
    $\xi = 1/6$ (conformal coupling) is preferred by the biquaternion
    structure is open.

  \item \textbf{[O] Off-shell consistency}: The Θ-only closure of the
    master equation (Step 6 of the GR chain) remains an open gap.

  \item \textbf{[O] Comparison with SM spectrum}: Systematic identification
    of UBT mode masses with Standard Model particle masses has not been
    completed.  See \texttt{research/fermion\_mass\_generation.tex}.
\end{enumerate}

% ======================================================================
\section{Cross-References}
\label{sec:xref}
% ======================================================================

\begin{tabular}{ll}
\hline
Topic & File \\
\hline
Canonical action & \texttt{consolidation\_project/appendix\_A\_theta\_action.tex} \\
Quantum field theory of $\Theta$ & \texttt{research/theta\_quantum\_field.tex} \\
GR chain & \texttt{canonical/bridges/GR\_chain\_bridge.tex} \\
Gauge group emergence & \texttt{canonical/bridges/gauge\_emergence\_bridge.tex} \\
Fermion masses & \texttt{research/fermion\_mass\_generation.tex} \\
Three generations & \texttt{research/generation\_structure.tex} \\
\hline
\end{tabular}

\end{document}
